\section{Methodology}
\label{sec:methodology}
\subsection{Study sites}

Experiments are conducted at two meteorological stations in Colombia that represent
climatologically contrasting regimes (Table~\ref{tab:sites}).

\begin{table}[t]
  \centering
  \caption{Study site metadata.}
  \label{tab:sites}
  \begin{tabular}{lcc}
    \toprule
    & \textbf{El Paso (César)} & \textbf{Uniandes (Bogotá)} \\
    \midrule
    Latitude  & $9.737^{\circ}\text{N}$ & $4.602^{\circ}\text{N}$ \\
    Longitude & $73.695^{\circ}\text{W}$ & $74.066^{\circ}\text{W}$ \\
    Elevation & 50~m a.s.l.      & 2\,600~m a.s.l.  \\
    Climate   & Tropical semi-arid & Equatorial mountain \\
    Avg.\ GHI & ${\sim}$5.5~kWh/m²/day & ${\sim}$4.2~kWh/m²/day \\
    \bottomrule
  \end{tabular}
\end{table}

\textbf{El Paso} is located in the lowland Caribbean region of C{\'e}sar department,
characterised by high annual irradiance, predominantly clear-sky conditions, and
dry-season cloud cover driven by the Inter-Tropical Convergence Zone (ITCZ).

\textbf{Universidad de los Andes} (referred to as Uniandes) is an instrumented
rooftop station in the Colombian Andean city of Bogot{\'a}. At 2\,600~m above sea
level, it experiences the equatorial bimodal rainfall regime with intense convective
cloud formation typically peaking in the early afternoon, producing sharp intraday
GHI drops.

\subsection{Data}

\subsubsection{Ground measurements}

GHI is measured at 10-minute resolution and stored in time-zone-aware (UTC) Parquet
files. Missing values account for less than 2\% of the record at both sites and are
excluded during dataset construction.

\subsubsection{Satellite data and patch extraction}
\label{sssec:satellite}

We use the GOES-16 ABI Level-2 Multi-Band Cloud and Moisture Imagery Product
(ABI-L2-MCMIPF), which provides all 16 ABI spectral channels at a unified
10-minute cadence and ${\sim}2$~km spatial resolution in the visible bands.
Raw hourly ABI-L2-MCMIPF granules are pre-processed into compact per-hour
archives: each archive is a compressed NumPy array of shape
$(6, 16, P, P)$, holding the six 10-minute slots within that hour, the
$C = 16$ spectral channels, and a $P \times P = 16 \times 16$ pixel patch
($P = 16$) centred on the target site, in units consistent with each band's
native reflectance or brightness temperature. At inference/training time, a
single timestamp is resolved to its $(C, P, P) = (16, 16, 16)$ frame by
selecting the corresponding hourly archive and 10-minute slot, so no
redundant decoding of the original satellite granules is required.

\subsubsection{Manifests, sequence construction, and temporal split}
\label{sssec:manifest}

For the deep-learning models (ResNet-LSTM, GraphSAGE-LSTM, FlatMLP), forecast
samples are constructed as pairs $(X, y)$ where the input
$X \in \R^{L \times C \times P \times P}$ is a sequence of $L = 24$ satellite
patch frames covering the most recent 4 hours of observation (10-minute
cadence), and the target $y \in \R$ is the ground-measured GHI at the site
$H$ steps ahead ($H \in \{6, 18, 36\}$ steps, corresponding to 1, 3, and
6 hours).

A pre-computed \textbf{manifest} (one Parquet table per site and horizon)
enumerates every valid forecast sample: for each candidate target timestamp it
stores the label $y$, the list of $L = 24$ history timestamps
\texttt{history\_ts}, and a coverage flag confirming that (i) all $L$ history
steps exist in the ground-truth series and (ii) a satellite archive is
available for each of them. Samples for which any of the $L$ history frames
is missing on disk are filtered out at load time, so every batch the models
see contains a complete 4-hour satellite history. Two thin
\texttt{Dataset} wrappers consume the same manifest and patch archives but
reshape the retrieved frames differently for the two model families:

\begin{itemize}
  \item \textbf{\texttt{PatchSeqDataset}} (used by ResNet-LSTM and FlatMLP)
        returns each sample as an image-like sequence
        $X \in \R^{L \times C \times P \times P}$ --- i.e., $L$ stacked
        $(16, 16, 16)$ tensors preserving the 2-D pixel grid, ready to be
        consumed by 2-D convolutions.
  \item \textbf{\texttt{GraphSeqDataset}} (used by GraphSAGE-LSTM) returns
        each sample as a node-feature sequence $X \in \R^{L \times N \times C}$
        with $N = P^2 = 256$: every frame is transposed and reshaped from
        $(C, P, P)$ to $(P^2, C)$, so that each of the 256 pixels becomes a
        graph node carrying a $C = 16$-dimensional feature vector (its
        16-channel spectral signature at that location).
\end{itemize}

To prevent data leakage, train, validation, and test sets are defined by
non-overlapping calendar blocks (Table~\ref{tab:splits}), and the manifest
coverage check, splitting boundaries, and patch archives are shared across
all three deep-learning models, ensuring that ResNet-LSTM, GraphSAGE-LSTM,
and FlatMLP are trained and evaluated on \emph{exactly} the same samples.

\begin{table}[t]
  \centering
  \caption{Temporal data splits.}
  \label{tab:splits}
  \setlength{\tabcolsep}{4pt}
  \begin{tabular}{llccc}
    \toprule
    \textbf{Site} & \textbf{Split} & \textbf{Start} & \textbf{End} & $n$ (approx.) \\
    \midrule
    \multirow{3}{*}{El Paso}
      & Train & 2022-03-01 & 2023-06-30 & 69\,984 \\
      & Val   & 2023-07-01 & 2023-10-31 & 17\,568 \\
      & Test  & 2023-11-01 & 2024-03-07 & 18\,288 \\
    \midrule
    \multirow{3}{*}{Uniandes}
      & Train & 2023-09-01 & 2024-09-30 & 56\,850 \\
      & Val   & 2024-10-01 & 2024-12-31 & 13\,104 \\
      & Test  & 2025-01-01 & 2025-03-28 & 12\,384 \\
    \bottomrule
  \end{tabular}
\end{table}

\subsection{Models}

\subsubsection{Persistence baseline}

The persistence model assumes irradiance is unchanged over the forecast horizon:
$\hat{y}(t+H) = \text{GHI}(t)$. It is optimal under no-change dynamics and
serves as the standard benchmark for short-term solar forecasting.

\subsubsection{ResNet-LSTM: convolutional spatial encoder}
\label{ssec:resnet}

ResNet-LSTM is an encoder--decoder that treats each satellite frame as a
conventional multi-channel image and lets a convolutional network learn its
spatial representation, before a recurrent network reasons over the sequence
of representations.

\textbf{Spatial encoder (per frame).} Each frame $X_\ell \in \R^{C \times P
\times P}$ ($\ell = 1,\ldots,L$) is passed independently through a lightweight
ResNet \cite{he2016}: a $3\times3$ convolutional stem, followed by three
residual stages of \texttt{BasicBlock} units (two $3\times3$ convolutions with
batch normalisation and a skip connection). The first stage preserves the
$16\times16$ spatial resolution, while the second and third stages apply
stride-2 downsampling, taking the spatial grid through
$16 \to 8 \to 4$. A global-average-pooling layer then collapses the
$4\times4$ feature map to a single vector, which a final linear layer projects
to a $d_e$-dimensional frame embedding $\mathbf{e}_\ell \in \R^{d_e}$
(\texttt{emb\_dim}, tuned in $\{64,128,192,256\}$). The encoder's channel
width is controlled by \texttt{base}~$\in\{16,24,32,48\}$ (the three stages
use \texttt{base}, $2\times$\texttt{base}, and $4\times$\texttt{base}
channels respectively); the resulting parameter count ranges from
approximately 50k to 500k depending on the configuration. Because the same
encoder weights are applied to every one of the $L = 24$ frames, the model is
translation-equivariant and shares a single spatial template across the whole
4-hour history.

\textbf{Temporal encoder.} The $L$ frame embeddings are normalised
(\texttt{LayerNorm}) and stacked into a sequence
$(\mathbf{e}_1,\ldots,\mathbf{e}_L) \in \R^{L \times d_e}$, which a
single- or two-layer LSTM \cite{hochreiter1997} (\texttt{hidden\_t}, \texttt{n\_lstm\_layers})
processes step by step; only its final hidden state $\mathbf{h}_L \in
\R^{d_t}$ is retained. A two-layer MLP head
(\texttt{Linear} $\to$ \texttt{ReLU} $\to$ \texttt{Dropout} $\to$
\texttt{Linear}) maps $\mathbf{h}_L$ to the scalar point forecast $\hat{y}$.

\subsubsection{GraphSAGE-LSTM: graph-based spatial encoder}
\label{ssec:gsage}

GraphSAGE-LSTM shares the same encoder--decoder skeleton as ResNet-LSTM ---
spatial encoder per frame, then LSTM over the embedding sequence, then MLP
head --- but \emph{replaces the convolutional encoder with a graph neural
network} that treats individual pixels, rather than the whole frame, as the
basic spatial unit.

\textbf{From pixels to a graph.} Each $(C, P, P)$ frame is reshaped (by
\texttt{GraphSeqDataset}, Section~\ref{sssec:manifest}) into $N = P^2 = 256$
nodes, one per pixel, each carrying its native $C = 16$-channel spectral
vector as a node feature --- no convolution or pooling is applied beforehand.
A single $k$-nearest-neighbour graph $G = (V, E, \mathbf{w})$ is built once
over the $16\times16$ pixel grid (by Euclidean distance in 2-D pixel
coordinates) and reused for every frame and every sample: each node connects
to its $k$ closest pixels ($k \in \{4,8,12,16\}$, tuned per
site/horizon, Section~\ref{ssec:optuna}), and every edge is assigned an
inverse-distance weight $w_{vu} = 1/d(v,u)$ (Eq.~\ref{eq:edge_weight} below).
Larger $k$ values progressively
add diagonal, two-step, and knight-move neighbours, so $k$ effectively
controls the spatial receptive field of a single graph layer. Unlike a fixed
8-connected grid or a convolutional filter, this $k$-NN topology lets the
model learn \emph{different} aggregation weights for differently-located
pixels (e.g.\ patch centre vs.\ periphery), rather than sliding one shared
template across the image. The edge index and weights are precomputed once,
registered as model buffers, and replicated across the batch dimension, so
the same mini-batch \texttt{DataLoader} and training loop used for
ResNet-LSTM apply unchanged.

\begin{equation}
  w_{vu} = \frac{1}{d(v,u)}, \quad (v, u) \in E.
  \label{eq:edge_weight}
\end{equation}

\textbf{Spatial encoder (per frame).} GraphSAGE \cite{hamilton2017} updates
every node's representation by aggregating its neighbourhood. With edge
weights, the neighbourhood mean becomes a distance-weighted average, and the
update concatenates each node's own (self) representation with that
weighted-neighbour summary before a shared linear projection and ReLU:
\begin{equation}
  \mathbf{h}_v^{(k)} = \sigma\!\left(
    \mathbf{W}^{(k)} \left[ \mathbf{h}_v^{(k-1)} \Vert
    \frac{\sum_{u \in \mathcal{N}(v)} w_{vu}\,\mathbf{h}_u^{(k-1)}}
         {\sum_{u \in \mathcal{N}(v)} w_{vu}} \right]
  \right)
  \label{eq:sage}
\end{equation}
where $\mathcal{N}(v)$ denotes the $k$-NN neighbours of $v$, $\Vert$ is
concatenation, and $\sigma$ is ReLU. The same layer (with shared weights
$\mathbf{W}^{(k)}$, output width \texttt{hidden\_g}) is applied to all 256
nodes simultaneously and stacked $K \in \{2,3,4\}$ times
(\texttt{n\_sage\_layers}); each additional layer extends a node's effective
receptive field by one more $k$-NN hop. A final \emph{mean readout} over all
$N = 256$ nodes collapses the per-pixel representations into a single
graph-level (frame) embedding $\mathbf{g}_\ell \in \R^{d_g}$, with
$d_g = $~\texttt{hidden\_g}~$\in \{64,96,128,192,256\}$ --- the GraphSAGE
analogue of the ResNet's $d_e$-dimensional frame embedding.

\textbf{Temporal encoder.} As in ResNet-LSTM, the $L$ graph-level embeddings
are normalised and stacked into $(\mathbf{g}_1,\ldots,\mathbf{g}_L) \in
\R^{L \times d_g}$ and processed by a single- or two-layer LSTM whose final
hidden state feeds an identical MLP head producing $\hat{y}$. The two
architectures are thus directly comparable: they differ \emph{only} in how
the per-frame spatial embedding is produced --- convolution-and-downsampling
over a regular grid (ResNet) versus weighted message-passing over a learned
$k$-NN pixel graph at full $256$-node resolution, with no spatial downsampling
(GraphSAGE) --- while sharing the same temporal model, output head, training
protocol, and data pipeline.

\subsubsection{FlatMLP: temporal-only ablation}
\label{ssec:mlp}

To isolate the contribution of \emph{spatial} encoding from that of temporal
modelling, FlatMLP discards all pixel-level spatial structure. For every
frame $X_\ell \in \R^{C \times P \times P}$, a global spatial average pool
collapses the $P \times P$ grid to its per-channel mean, $\bar{\mathbf{x}}_\ell
\in \R^{C}$ --- i.e.\ each frame is reduced to a 16-dimensional ``patch-mean
spectral signature'' with no notion of pixel location. The $L = 24$ resulting
vectors are concatenated into a single flat input
$\bar{\mathbf{x}} \in \R^{L \times C} = \R^{384}$ ($24 \times 16$). This
vector is passed through a fully-connected feed-forward network of
\texttt{n\_layers} blocks, each consisting of \texttt{Linear} $\to$
\texttt{LayerNorm} $\to$ \texttt{ReLU} $\to$ \texttt{Dropout}
(\texttt{hidden\_dim} units), followed by a final \texttt{Linear} layer that
outputs the scalar $\hat{y}$ directly --- every unit in one layer is connected
to every unit in the next (a standard dense MLP), with no recurrence, no
convolution, and no graph structure. Because FlatMLP receives exactly the
same $L\times C \times P \times P$ input tensors (via \texttt{PatchSeqDataset})
and is trained with the same protocol, its skill gap relative to ResNet-LSTM
and GraphSAGE-LSTM directly quantifies how much of their forecasting skill
stems from \emph{spatial} encoding of the satellite patch, as opposed to
simple per-frame statistics processed over time.

\subsection{Hyperparameter optimisation}
\label{ssec:optuna}

For the three deep-learning models (ResNet-LSTM, GraphSAGE-LSTM, and FlatMLP)
we use Optuna \cite{akiba2019} with the TPE sampler and Median pruner.
The study budget is 50 trials per configuration
(site $\times$ horizon $\times$ seed) with $n_\text{jobs}=2$; an extended
run with 100 trials is being evaluated to confirm convergence.

The ResNet-LSTM and GraphSAGE-LSTM search spaces (Table~\ref{tab:hpspace})
include \texttt{n\_lstm\_layers}~$\in\{1,2\}$ to allow stacked recurrent
processing, and a categorical L1 regularisation coefficient
(\texttt{l1\_reg}) to penalise large weights in high-capacity configurations;
both are new relative to the initial baseline study.
For \textbf{ResNet-LSTM} specifically, the embedding dimension
(\texttt{emb\_dim}) and temporal hidden size (\texttt{hidden\_t}) are extended
to 256 to increase capacity.
For \textbf{GraphSAGE-LSTM}, the number of nearest neighbours
($k \in \{4, 8, 12, 16\}$) used to build the weighted pixel graph is
additionally tunable, allowing the study to determine the optimal spatial
connectivity for each site-horizon combination.
For \textbf{FlatMLP}, the study searches over MLP depth, hidden dimension,
dropout, learning rate, and weight decay (50 trials).

\begin{table*}[t]
  \centering
  \caption{Optuna hyperparameter search spaces. L1 coefficient and LSTM depth are
           new relative to the initial baseline study. The \texttt{k\_neighbors}
           parameter determines the weighted $k$-NN graph built per trial for
           GraphSAGE-LSTM.}
  \label{tab:hpspace}
  \setlength{\tabcolsep}{5pt}
  \begin{tabular}{lll}
    \toprule
    \textbf{Hyperparameter} & \textbf{ResNet-LSTM} & \textbf{GraphSAGE-LSTM} \\
    \midrule
    \texttt{base} / \texttt{hidden\_g}           & $\{16,24,32,48\}$                     & $\{64,96,128,192,256\}$               \\
    \texttt{emb\_dim} / \texttt{n\_sage\_layers} & $\{64,128,192,256\}$                  & $\{2,3,4\}$                           \\
    \texttt{hidden\_t} (LSTM)                    & $\{64,128,192,256\}$                  & $\{64,96,128,192,256\}$               \\
    \texttt{n\_lstm\_layers}                     & $\{1,2\}$                             & $\{1,2\}$                             \\
    \texttt{k\_neighbors}                        & ---                                   & $\{4,8,12,16\}$                       \\
    \texttt{l1\_reg}                             & $\{0,\,10^{-5},\,10^{-4},\,10^{-3}\}$& $\{0,\,10^{-5},\,10^{-4},\,10^{-3}\}$\\
    \texttt{dropout}                             & $\{0,0.1,0.2,0.3\}$                  & $\{0,0.1,0.2\}$                       \\
    \texttt{lr}                                  & $[3{\times}10^{-4},\;3{\times}10^{-3}]$ & $[5{\times}10^{-5},\;2{\times}10^{-3}]$ \\
    \texttt{weight\_decay}                       & $[10^{-6},\;10^{-3}]$                & $[10^{-6},\;10^{-3}]$                 \\
    \texttt{batch\_size}                         & $\{16,32,64\}$                        & $\{8,16,32\}$                         \\
    \bottomrule
  \end{tabular}
\end{table*}

The Optuna study identifies the best trial on validation $\rmseday$; the winning
configuration is then retrained for up to 30~epochs with patience~8 and evaluated
on the held-out test set.

\subsection{Training protocol}

All deep-learning models are trained with the Adam optimiser minimising the MSE
loss on normalised targets ($y_\text{norm} = (y - \bar{y}) / s_y$, with $\bar{y}$
and $s_y$ estimated from the training set). Training uses early stopping with
patience~6 (trials) or~8 (final retraining). Automatic Mixed Precision (AMP) is
enabled when a GPU is available. Each experiment is replicated with four seeds
$\{42, 1, 7, 13\}$ to report mean~$\pm$~std statistics.

\subsection{Uncertainty quantification: Split Conformal Prediction}
\label{ssec:uq}

All point-forecast models are used as the backbone for a three-stage uncertainty
quantification (UQ) pipeline whose long-term goal is to decompose total forecast
uncertainty into its \emph{aleatoric} component (irreducible cloud stochasticity)
and its \emph{epistemic} component (uncertainty due to limited data and model
capacity).

\textbf{Stage~1 --- Split Conformal Prediction (implemented, \cite{vovk2005,angelopoulos2023}).}
We apply Inductive Conformal Prediction as a distribution-free, finite-sample
calibration method.  A non-conformity score is computed over the validation set:
\begin{equation}
  s_i = |y_i - \hat{y}_i| \quad [\text{W/m}^2].
  \label{eq:score}
\end{equation}
The calibrated conformal quantile at miscoverage level $\alpha$ is
\begin{equation}
  \hat{q} = \mathrm{Quantile}\!\left(\{s_i\}_{i=1}^{n_\text{cal}},\;
    \frac{\lceil(n_\text{cal}+1)(1-\alpha)\rceil}{n_\text{cal}}\right),
  \label{eq:qhat}
\end{equation}
and the prediction interval for a new test point is
\begin{equation}
  \mathcal{C}_{1-\alpha}(x) = \bigl[\hat{y}(t+H)-\hat{q},\;\hat{y}(t+H)+\hat{q}\bigr].
  \label{eq:interval}
\end{equation}
Under exchangeability, this guarantees marginal coverage
$\mathbb{P}(Y \in \mathcal{C}_{1-\alpha}) \geq 1-\alpha$ exactly in finite samples
\cite{angelopoulos2023}.  Because $\hat{q}$ is a global constant, the interval
has fixed width $2\hat{q}$ and does not adapt to local prediction difficulty.
We evaluate empirical coverage at $\alpha \in \{0.05, 0.10, 0.20\}$ (intervals
of 95\%, 90\%, and 80\%) reporting both overall and daytime-only metrics.

\textbf{Stage~2 --- Conformalized Quantile Regression (planned, \cite{romano2019}).}
To obtain intervals whose width adapts to local uncertainty, we will train
quantile-regression heads on the DL encoders and apply the CQR correction.
This is expected to reduce unnecessary interval width during low-variability
periods (clear-sky mornings) without sacrificing coverage.

\textbf{Stage~3 --- Bayesian uncertainty decomposition (planned,
\cite{kendall2017,welling2011}).}
Gaussian variance heads (NLL loss) will model aleatoric uncertainty, while
Stochastic Gradient Langevin Dynamics (SGLD, \cite{welling2011}) will sample
the posterior weight distribution to quantify epistemic uncertainty.
SGLD injects calibrated Langevin noise into gradient updates, enabling
approximate Bayesian inference without the computational overhead of full
MCMC: running the chain for $T$ epochs and collecting snapshots every
$\Delta$ epochs yields $T/\Delta$ approximate posterior samples, whose
ensemble mean serves as the point forecast and whose standard deviation
quantifies model-knowledge uncertainty.
Together, the aleatoric head and SGLD ensemble will decompose total forecast
uncertainty into its irreducible (cloud stochasticity) and reducible
(model-knowledge) components \cite{kendall2017}, supporting risk-stratified
dispatch decisions for solar-integrated power systems.

\subsection{Evaluation metrics}

Following standard practice in solar forecasting \cite{coimbra2013}, we report:

\textbf{Daytime RMSE} ($\rmseday$, W/m²): RMSE computed only over test samples
with observed $\text{GHI} \geq 20$~W/m², excluding night-time samples where all
models trivially predict zero.

\textbf{Daytime skill score} ($\skillday$): normalised improvement over persistence,
\begin{equation}
  \skillday = 1 - \frac{\rmseday(\text{model})}{\rmseday(\text{persistence})}
  \label{eq:skill}
\end{equation}
A score of 0 indicates parity with persistence; positive values indicate improvement.
We also report $\text{MAE}_\text{day}$ for completeness.